\section{ماژول liftr (Linearized Function Trick Compressor)}
\label{sec:module_liftr}

\subsection{معرفی و انگیزه}
این مکانیزم بر مبنای مبانی نظری پیشرفته در ادبیات توجه و فشرده‌سازی \cite{ref_liftr_2024} توسعه یافته است.

پیچیدگی محاسباتی درجه دوم $\mathcal{O}(N^2)$ در لایه‌های توجه استاندارد مبتنی بر بیشینه‌ساز هموار (\lr{Softmax Attention}) و همچنین محاسبات غیرخطی توابع فعال‌سازی نظیر \lr{GELU} و \lr{SiLU}، گلوگاه اصلی سرعت در پردازش بافت‌های طولانی به شمار می‌آیند. راهکارهای خطی‌سازی معمولاً نیازمند بازآموزی کامل شبکه از ابتدا با توابع هسته پیچیده هستند که هزینه‌های گزافی را تحمیل می‌کند.

ماژول \lr{liftr} تحت عنوان ترفند خطی‌سازی توابع (\lr{Linearized Function Trick}) برای حذف فوری غیرخطی‌های محاسباتی بدون نیاز به آموزش سنگین اولیه توسعه یافته است. در این ماژول دو استراتژی مکمل به کار گرفته می‌شود: اول، جایگزینی تابع \lr{Softmax} در مکانیزم توجه با میانگین‌گیری سرتاسری تانسور ارزش‌ها (\lr{Mean-Value Projection}) که پیچیدگی را مستقیماً به $\mathcal{O}(N)$ کاهش می‌دهد؛ و دوم، خطی‌سازی توابع فعال‌سازی غیرخطی در لایه‌های پیش‌خوران به صورت نگاشت همانی مقیاس‌گذاری‌شده ($f(x) \approx \alpha \cdot x$).

\subsection{مبانی نظری و فرمول‌بندی ریاضی}
در لایه توجه نرم استاندارد، خروجی به صورت زیر تعریف می‌شود:
\begin{equation}
\text{Attention}(\mathbf{Q}, \mathbf{K}, \mathbf{V}) = \text{Softmax}\left( \frac{\mathbf{Q}\mathbf{K}^\top}{\sqrt{d}} \right) \mathbf{V}
\end{equation}
روش \lr{LIFTR} این فرآیند را با تقریب مرتبه صفر ماتریس وزن‌های توجه جایگزین می‌نماید که در آن توزیع احتمال روی توالی با توزیع یکنواخت تقریب زده می‌شود:
\begin{equation}
\text{Softmax}\left( \frac{\mathbf{Q}\mathbf{K}^\top}{\sqrt{d}} \right)_{ij} \approx \frac{1}{N}
\end{equation}
در نتیجه، خروجی زمینه توجه به صورت بردار میانگین توالی ارزش‌ها ساده‌سازی می‌شود:
\begin{equation}
\mathbf{C} = \frac{1}{N} \sum_{i=1}^N \mathbf{V}_i \in \mathbb{R}^{B \times 1 \times d_v}
\end{equation}
این بردار سپس در راستای بعد توالی بسط یافته (\lr{Broadcast/Expand}) و از طریق ماتریس پروجکشن خروجی پرتاب می‌گردد:
\begin{equation}
\mathbf{Y} = \text{expand}(\mathbf{C}, [B, N, d_v]) \, \mathbf{W}_o + \mathbf{b}
\end{equation}
این فرآیند نیاز به ذخیره ماتریس ضرب داخلی $N \times N$ و عملگر نمایی \lr{Softmax} را به کلی ملغی می‌سازد.

همچنین برای فعال‌سازی‌های غیرخطی \lr{GELU} و \lr{SiLU}، تقریب درجه اول در همسایگی دامنه کاری متداول اعمال می‌شود:
\begin{equation}
\text{GELU}(x) \approx \alpha \cdot x, \quad \alpha \in [0.7, 1.0]
\end{equation}
که در پیاده‌سازی مرجع مقدار ضریب مقیاس به صورت پیش‌فرض $\alpha = 0.8$ تنظیم گردیده است.

\paragraph{تحلیل طیفی و نامساوی‌های پایداری خطا}
برای بررسی میزان دقت تقریب ماتریس وزن یا نگاشت توجه در ماژول \lr{liftr}، تفاضل ماتریس فاکتورگیری‌شده از نگاشت متراکم اصلی بر اساس نامساوی هافدینگ-فون‌نیومان ارزیابی می‌گردد. فرض کنید $\mathbf{W} \in \mathbb{R}^{d_{\text{in}} \times d_{\text{out}}}$ ماتریس وزن هدف باشد که توسط تقریب کم‌مرتبه یا پراکنده با رتبه موثر $r$ جایگزین شده است. کران بالای خطای بازسازی بر مبنای نرم فروبنیوس به صورت زیر تصریح می‌شود:
\begin{equation}
\label{eq:spectral_bound_liftr}
\| \mathbf{W} - \mathbf{W}_{r} \|_F = \sqrt{\sum_{i=r+1}^{\min(d_{\text{in}}, d_{\text{out}})} \sigma_i^2(\mathbf{W})} \le \epsilon \| \mathbf{W} \|_F
\end{equation}
که در آن $\sigma_i(\mathbf{W})$ بیانگر $i$-امین مقدار منفرد مرتب‌شده ماتریس به صورت نزولی است. نسبت حفظ انرژی طیفی $\eta_r = \frac{\sum_{i=1}^r \sigma_i^2}{\sum_{i=1}^{d} \sigma_i^2}$ تضمین می‌کند که با انتخاب آستانه $\tau \ge 0.90$، بیش از نود درصد از پراکندگی واریانس داده‌ها حفظ شده و افت سیگنال در انتشار گرادیان رو به عقب به حداقل ممکن کاهش می‌یابد.

\paragraph{تحلیل مجانبی و پیچیدگی اشتراک حافظه سخت‌افزاری}
از منظر پیچیدگی محاسباتی و انتقال داده‌ها در سلسومراتب حافظه (حافظه سراسری \lr{HBM} به حافظه محلی \lr{SRAM})، عملیات ضرب ماتریسی در این مکانیزم نیازمند تخصیص بافر موقت با ابعاد فاکتورگیری‌شده است. برای پردازش دنباله‌ای به طول $N$ با بعد مخفی $D$، بار محاسباتی هسته از فرم درجه دو به رابطه زیر تقلیل می‌یابد:
\begin{equation}
\label{eq:complexity_liftr}
\mathcal{C}_{\text{liftr}} = \mathcal{O}\left( B \cdot N \cdot D \cdot r + B \cdot H \cdot N \cdot d_k \right) \ll \mathcal{O}\left( B \cdot H \cdot N^2 \cdot d_k \right)
\end{equation}
این بهبود مرتبه پیچیدگی، با کاهش چشمگیر حجم تبادل تانسورها در گذرگاه‌های ارتباطی کارت گرافیک، مانع از اشباع پهنای باند حافظه (پدیده \lr{Memory-Bound}) در فرآیند استنتاج طولانی‌مدت می‌گردد.

\subsection{معماری و پیاده‌سازی}
ماژول در فایل \lr{Codes/framework/compression\_engine/liftr.py} پیاده‌سازی شده و ساختار نرم‌افزاری آن به قرار زیر است:
\begin{itemize}
    \item \textbf{کلاس \lr{\_LinearizedGELU}:} جایگزین سبک لایه‌های فعال‌سازی \lr{nn.GELU} و \lr{nn.SiLU} که عملگر غیرخطی را با ضرب اسکالر ساده $\alpha \cdot x$ تعویض می‌کند.
    \item \textbf{کلاس \lr{\_LinearizedAttentionWrapper}:} لایه پوشاننده ساختاری توجه که با پشتیبانی هوشمند از انواع معماری‌ها، خروجی ارزش‌ها را استخراج و میانگین‌گیری می‌کند:
    \begin{itemize}
        \item پشتیبانی از خانواده مدل‌های مدرن (\lr{LLaMA/Mistral/Qwen}) از طریق لایه‌های \lr{v\_proj} و \lr{o\_proj}.
        \item پشتیبانی از انکودرهای \lr{DistilBERT} از طریق \lr{v\_lin} و \lr{out\_lin}.
        \item پشتیبانی از انکودرهای استاندارد \lr{RoBERTa/BERT} از طریق ماژول‌های \lr{value} و \lr{query}.
        \item پشتیبانی از دیکودرهای علّی \lr{GPT-2} از طریق تفکیک مؤلفه ارزش از لایه ادغامی \lr{c\_attn} و ارسال به \lr{c\_proj}.
    \end{itemize}
    \item \textbf{مکانیزم تفویض و بازگشت ایمن (\lr{Safe Fallback Delegation}):} در صورتی که معماری لایه توجه ناشناخته باشد یا شکل ورودی با شرایط ماژول منطبق نباشد، محاسبات بدون بروز خطا به لایه توجه اصلی هدایت می‌شود.
    \item \textbf{کلاس \lr{LIFTRCompressor}:} پیاده‌سازی موتور فشرده‌سازی جهت پیمایش گراف مدل و جایگزینی لایه‌ها.
\end{itemize}

\subsection{نتایج تجربی}
ارزیابی ماژول \lr{liftr} بر روی مدل‌های زبانی و انکودرهای پایه بر اساس داده‌های تجربی تاییدشده در \lr{thesis\_final/master\_verified\_results.json} انجام شده و در جدول \ref{tab:benchmark_liftr} گزارش شده است:

\begin{table}[H]
\centering
\small
\caption{نتایج تجربی و اعتبارسنجی بنچ‌مارک فشرده‌ساز \lr{liftr} روی مدل‌های پایه}
\label{tab:benchmark_liftr}
\begin{tabular}{lccccc}
\toprule
\textbf{مدل پایه} \lr{(Base Model)} & \textbf{دقت} \lr{(Accuracy)} & \textbf{شاخص} \lr{(F1-Score)} & \textbf{حافظه} \lr{(Peak VRAM)} & \textbf{توان} \lr{(Throughput)} & \textbf{وضعیت} \\
\midrule
\lr{DistilBERT} & \lr{90.14\%} & \lr{0.9029} & \lr{337.59 MB} & \lr{182,296.25 tok/s} & \lr{PASS} \\
\lr{RoBERTa} & \lr{80.16\%} & \lr{0.8309} & \lr{585.62 MB} & \lr{92,788.29 tok/s} & \lr{PASS} \\
\lr{GPT-2} & \lr{61.01\%} & \lr{0.6654} & \lr{634.91 MB} & \lr{68,286.26 tok/s} & \lr{PASS} \\
\lr{Qwen} & \lr{0.00\%} & \lr{---} & \lr{---} & \lr{---} & \lr{FAIL} \\
\bottomrule
\end{tabular}
\end{table}

داده‌های تجربی ثبت‌شده نشان‌دهنده توان عملیاتی قابل توجه این ماژول است. ماژول \lr{liftr} با حذف محاسبات سنگین توجه موفق شده است در انکودر \lr{DistilBERT} به بالاترین نرخ توان پردازشی در میان ماژول‌های مورد بررسی یعنی \lr{182,296.25 tok/s} با دقت درخشان \lr{90.14\%} دست یابد. در مدل‌های \lr{RoBERTa} و \lr{GPT-2} نیز به ترتیب دقت‌های \lr{80.16\%} و \lr{61.01\%} حاصل شده که نشان‌دهنده حفظ مناسب اطلاعات با وجود خطی‌سازی کامل است. در ارزیابی مدل \lr{Qwen} به علت شکست در تطابق ابعاد سرهای توجه گروهی (\lr{GQA}) در آزمون مستقیم خطوط لوله لایه‌های ادغامی، وضعیت \lr{FAIL} به ثبت رسیده است.

\subsection{محدودیت‌ها و تحلیل شکست}
تحلیل مکانیزم‌های خطی‌سازی ماژول \lr{liftr} محدودیت‌های بنیادی زیر را نمایان می‌سازد:
\begin{itemize}
    \item \textbf{فقدان فشرده‌سازی فیزیکی پارامترها (\lr{No Parameter Reduction}):} ماژول \lr{liftr} ماهیت فشرده‌سازی محاسباتی (\lr{Compute Complexity Reduction}) دارد و ماتریس‌های وزن را حذف یا هرس نمی‌کند؛ بنابراین حافظه ایستا اشغال‌شده توسط وزن‌های مدل بدون تغییر باقی می‌ماند.
    \item \textbf{از بین رفتن وابستگی‌های مکانی توکن‌ها (\lr{Loss of Spatial Resolution}):} تبدیل ماتریس توجه به میانگین یکنواخت مقادیر ارزش‌ها موجب می‌شود که اطلاعات مربوط به ترتیب و وزن‌دهی اختصاصی میان جفت توکن‌ها از بین برود؛ این امر در تسک‌هایی که نیازمند استدلال مکانی یا ردیابی دقیق مؤلفه‌ها در بافت‌های پیچیده هستند موجب افت کیفیت بازنمایی می‌گردد.
    \item \textbf{ناسازگاری با سرهای توجه گروهی در مدل \lr{Qwen} (\lr{GQA Incompatibility}):} عدم تطابق نسبت تعداد سرهای کوئری به کلید/ارزش در ساختار \lr{Grouped-Query Attention} منجر به بروز خطا در عملگر پخش (\lr{Broadcast}) ابعاد میانگین ارزش‌ها شده که عامل ریشه‌ای شکست ارزیابی در مدل \lr{Qwen} بوده است.
\end{itemize}
